\chapter{Trabajo Futuro} \label{sec:TrabajosFuturos}

El marco de pruebas desarrollado puede ampliarse para mitigar las limitaciones estructurales identificadas durante el diseño del entorno experimental.

\section{Implementación de Admission Control}
El flujo GitOps actual asegura la convergencia del estado declarado, pero no ejecuta validación semántica pre-despliegue. Un trabajo futuro inmediato es integrar controladores de admisión como OPA/Gatekeeper o Kyverno en el pipeline, lo que evitaría la inyección al clúster de políticas de red mal configuradas, redundantes o sintácticamente válidas pero lógicamente inseguras.

\section{Detección y Respuesta en Tiempo Real (Runtime Security)}
Las Network Policies demostraron prevenir conexiones no autorizadas, pero no proveen observabilidad sobre los intentos de ataque bloqueados. Se propone desplegar herramientas de seguridad en tiempo real basadas en eBPF, como Falco o Tetragon, para monitorizar las llamadas al sistema operativo y generar alertas activas frente a la violación sistemática de los límites de red.

\section{Integración de Identidad Criptográfica (Service Mesh)}
Las políticas evaluadas gestionan el aislamiento en capas L3/L4 (y L7 en Cilium) mediante resolución de etiquetas e IPs. Una extensión necesaria consiste en medir el overhead de combinar estos CNIs con un plano de control de Service Mesh (ej. Istio o Linkerd) para inyectar proxies sidecar y proveer autenticación mutua (mTLS), evaluando el costo de CPU y latencia del cifrado en el espacio de usuario frente al datapath nativo del CNI.

\section{Pruebas de Estrés en Mapas eBPF L7}
Si bien Cilium probó ser efectivo para el filtrado HTTP básico, la prueba se limitó a un entorno delimitado. Un trabajo derivado consiste en someter el plano de datos de Cilium a arquitecturas de red con miles de reglas granulares en capa 7 (filtrando rutas específicas, métodos gRPC y topics de Kafka simultáneamente) para establecer el umbral exacto donde el consumo de RAM del mapa eBPF empieza a degradar las métricas de rendimiento.
